% !TEX root = ../main.tex
\chapter{The Money View, Micro and Macro}\label{ch:moneyview}
\begin{paracol}{2}

\begin{sources}
\begin{tabularx}{\linewidth}{@{}l l L@{}}
\toprule
\textbf{Segment} & \textbf{Length} & \textbf{Content}\\
\midrule
\seg{4}{1}{L4.1 FT: Dealer of Last Resort} & 5:23 & Bernanke's QE3 and Draghi's conditional promise as dealer-of-last-resort actions.\\
\seg{4}{2}{L4.2 Reading: Hyman Minsky} & 3:07 & Minsky and the financial instability hypothesis; Copeland's flow of funds.\\
\seg{4}{3}{L4.3 Payments: Money and Credit} & 5:36 & Paying with money, with pure credit, with a bank's IOU.\\
\seg{4}{4}{L4.4 Payments: Discipline and Elasticity} & 4:13 & The payment system is a credit system.\\
\seg{4}{5}{L4.5 The Survival Constraint} & 3:35 & Every agent as a bank: cash inflow must cover cash commitments.\\
\seg{4}{6}{L4.6 Sources and Uses Accounts} & 6:55 & The sources-and-uses matrix and its two rules.\\
\seg{4}{7}{L4.7 Payment Example: Money and Credit} & 10:05 & A coffee paid in cash, a dinner paid by credit card.\\
\seg{4}{8}{L4.8 Flow of Funds Accounts} & 10:10 & Copeland's macro accounts; statistical discrepancies; what they miss.\\
\seg{4}{9}{L4.9 The Survival Constraint, redux} & 2:40 & Three ways to meet a cash shortfall.\\
\seg{4}{10}{L4.10 Liquidity, Long and Short} & 9:42 & Time pattern of cash flows and commitments; price-insensitive borrowers.\\
\seg{4}{11}{L4.11 Financial Fragility, Flows and Stocks} & 6:23 & The scales of cash flows and commitments; the money market as stress detector.\\
\bottomrule
\end{tabularx}

\smallskip
\textbf{Files.} Transcripts \file{materials/transcripts/L04-P*.txt}.
\textbf{Reading.} Mehrling's article on Hyman Minsky (not among the course files); Minsky, \emph{Stabilizing an Unstable
Economy} (1986), quoted in class.
\end{sources}

\section*{Motivation}
Lectures 1--3 drew the hierarchy with balance sheets, which are \emph{stocks}. This lecture introduces the other half of the
method: \emph{flows} of cash. Its central idea, taken from Hyman Minsky, is to look at every economic unit ``as if it were a
bank'', facing every day a \emph{survival constraint}: cash inflows must cover cash commitments. The lecture develops it at two
levels:
\begin{enumerate}
  \item \textbf{micro}: payments with money and with credit, the sources-and-uses accounts of a single agent, the survival
        constraint and the ways to meet it;
  \item \textbf{macro}: Morris Copeland's flow of funds accounts, and the money market interest rate as the place where the
        balance of the whole economy's cash flows and commitments shows up.
\end{enumerate}

% ------------------------------------------------------------------
\section{FT: dealer of last resort}\label{sec:mv-ft}

\begin{casestudy}[FT, September 2012: QE3 and Münchau on the ECB]
The FOMC has announced that it will keep buying mortgage-backed securities until employment improves: ``quite astonishing'' open-ended
support, prepared by Bernanke's remarks at Jackson Hole \ts{4}{1}{0:07}\ts{4}{1}{0:27}. In the same week Draghi said he would do
whatever is necessary: central bankers are determined to use their balance sheets to prevent anything bad \ts{4}{1}{0:48}. In the FT,
Wolfgang Münchau argues that QE would be the right policy for Europe too \ts{4}{1}{1:09}. He points out that Draghi will buy Spanish or
Italian bonds only if those countries \emph{ask}, which a political leader may be reluctant to do, while Bernanke simply announced
that he would buy \ts{4}{1}{1:30}\ts{4}{1}{1:50}\ts{4}{1}{2:11}.
\end{casestudy}

\paragraph{Promises and purchases.} If Spain and Italy do not ask, why should the price of their bonds move? Only through
expectations, and Münchau, with Mehrling, considers that a weak reed \ts{4}{1}{2:32}. The course looks at what central banks actually
do \ts{4}{1}{2:53}. Mehrling will call both actions \emph{dealer of last resort} actions: the central bank promises to buy bonds at a
particular price, so that no one will ever sell them for less; buying at a high price lowers their yield \ts{4}{1}{3:16}\ts{4}{1}{3:37}.
For this to work the promise must be credible, i.e.\ the central bank must actually buy. Bernanke is buying, and the market price will
move to his price at once; Mehrling predicts that the price of Spanish and Italian bonds will not move all the way to Draghi's level
until Draghi actually buys \ts{4}{1}{3:57}\ts{4}{1}{4:18}.

\begin{intuition}[Why central banks are so bold]
Central banks are doing things they have never done, and boldly, because the fiscal authorities are not: in Europe they are caught up in
austerity and in blaming Greece, in the US Republicans and Democrats cannot agree. The central banks are trying to hold the system
together through the election cycle until fiscal authorities can sort things out, ``not a job that any central banker ever signed up
for'' \ts{4}{1}{4:38}\ts{4}{1}{4:58}.
\end{intuition}\begin{coursenote}[OMT conditionality]
Outright Monetary Transactions (\cref{sec:nh-ft}) require that the country first request assistance from the European Stability
Mechanism and accept its conditions; this is the ``asking'' that Münchau worried about.
\end{coursenote}

% ------------------------------------------------------------------
\section{Reading: Hyman Minsky}\label{sec:mv-minsky}

Mehrling knew Hyman Minsky, who belongs to the American tradition he calls his own (Wesley Clair Mitchell, Edward Shaw, Morris Copeland)
\ts{4}{2}{0:07}. When the crisis started, the \emph{New Yorker} ran a ``Talk of the Town'' piece calling it a ``Minsky moment''; for a man
who grew up in New York, after decades of neglect by economists, that would have been immortality \ts{4}{2}{0:28}\ts{4}{2}{0:49}.

\begin{reading}[Mehrling on Minsky]
Minsky's \emph{financial instability hypothesis}: there is something inherently unstable about financial systems, and the role of the
central bank is to put bounds on that instability, to keep it from flying off on the upside or on the downside; this is not easy
\ts{4}{2}{1:10}. It is a very different vision from the technocratic one of moving the Fed funds rate by another 25 basis points to get
it just right; Minsky's was a mentality prepared for financial crisis \ts{4}{2}{1:30}\ts{4}{2}{1:51}. A common misreading, Mehrling warns,
takes Minsky's hedge, speculative and Ponzi units to be about leverage; read through Copeland, they are about the \emph{alignment in time}
of cash flows and commitments (\cref{sec:mv-longshort}) \ts{4}{10}{0:07}.
\end{reading}\begin{history}[The Minsky moment]
Hyman P.~Minsky (1919--1996), born in Chicago, taught at Washington University in St.~Louis and ended his career at the Levy Economics
Institute of Bard College. The phrase ``Minsky moment'' was coined by Paul McCulley in 1998 for the Russian crisis; John Cassidy's
\emph{New Yorker} piece ``The Minsky Moment'' appeared on 4 February 2008. Mehrling's ``grew up in New York'' refers to Minsky's
youth: he attended high school in New York City.\par
\emph{References:} H.~P.~Minsky, \emph{Stabilizing an Unstable Economy}, Yale UP (Twentieth Century Fund), 1986; J.~Cassidy,
\emph{The New Yorker}, 4 Feb 2008.
\end{history}

\paragraph{Three images of the macroeconomy.} Macroeconomics courses give two images of how money enters the economy
\ts{4}{2}{2:31}:
\begin{enumerate}
  \item the national income and product accounts (NIPA), $C+I+G+NX = Y$;
  \item the quantity theory, $MV = PT$.
\end{enumerate}
A third, proposed by Morris Copeland in 1952, is the \emph{flow of funds}: used on Wall Street and inside the Fed, but never turned into a
proper macroeconomic theory, which many think is what should be attempted \ts{4}{2}{2:52}.

% ------------------------------------------------------------------
\section{Payments: money and credit}\label{sec:mv-payments}

Suppose there are only two people, you and I; each produces what the other wants, but not at the same time, so barter does not work
\ts{4}{3}{0:07}\ts{4}{3}{0:28}. There are three ways to arrange the trade (I buy goods from you):
\begin{enumerate}
  \item \textbf{Money}: you give me goods, I give you a token; only existing money changes pocket \ts{4}{3}{0:56}\ts{4}{3}{1:18}.
  \item \textbf{Pure credit}: I write you a note, ``I owe you''; this requires trust, as between friends who take turns paying lunch
        \ts{4}{3}{1:39}\ts{4}{3}{2:00}\ts{4}{3}{2:41}.
  \item \textbf{Bank credit}: if we do not trust each other but both trust a third party, the bank accepts my IOU and issues its own IOU,
        which you accept \ts{4}{3}{3:42}\ts{4}{3}{4:02}\ts{4}{3}{4:24}.
\end{enumerate}
\begin{taccts}
\textit{(1) Money:}\quad
\tacct[1.9cm]{Me}{$+$Goods\\$-\Delta M$}{}\quad
\tacct[1.9cm]{You}{$-$Goods\\$+\Delta M$}{}

\medskip
\textit{(2) Pure credit:}\quad
\tacct[1.9cm]{Me}{$+$Goods}{$+\Delta$IOU}\quad
\tacct[1.9cm]{You}{$-$Goods\\$+\Delta$IOU}{}

\medskip
\textit{(3) Bank credit:}\quad
\tacct[1.6cm]{Me}{$+$Goods}{$+$IOU}\quad
\tacct[1.6cm]{Bank}{$+$IOU (mine)}{$+M$}\quad
\tacct[1.6cm]{You}{$-$Goods\\$+M$}{}
\end{taccts}
The difference between (1) and the other two: with money the quantity of money does not change, existing pieces of paper move; with
credit the quantity of outstanding IOUs \emph{increases} with the transaction \ts{4}{3}{3:01}\ts{4}{3}{3:21}. In (3) the bank's IOU is
called $M$ because it is used as money, but it is a liability of the bank created by expanding its balance sheet: money is a form of
credit, and it expands and contracts with the pattern of trade \ts{4}{3}{4:45}\ts{4}{3}{5:06}. This stripped-down picture is how the
modern economy works \ts{4}{3}{5:27}.

\section{Discipline and elasticity in payments}\label{sec:mv-de}

\begin{nfigure}
\begin{tikzpicture}[font=\scriptsize,x=0.55cm,y=0.55cm]
  \begin{scope}
    \draw[-{Stealth[length=3pt]}] (0,0) -- (9,0) node[right]{time};
    \draw[-{Stealth[length=3pt]}] (0,0) -- (0,5);
    \draw[dashed] (0,4) -- (8.5,4) node[right]{total $M$};
    \draw[thick,defcol] plot[smooth] coordinates {(0,2) (1,2.8) (2,3.6) (3,2.6) (4,1.2) (5,0.3) (6,1.5) (7,2.9) (8,2.2)};
    \draw[thick,thmcol] plot[smooth] coordinates {(0,2) (1,1.2) (2,0.4) (3,1.4) (4,2.8) (5,3.7) (6,2.5) (7,1.1) (8,1.8)};
    \node[defcol] at (8.8,2.4) {mine}; \node[thmcol] at (8.8,1.6) {yours};
    \node[font=\small\bfseries] at (4,5.5) {money: fixed total};
  \end{scope}
  \begin{scope}[xshift=7.4cm,yshift=1.4cm]
    \draw[-{Stealth[length=3pt]}] (0,0) -- (9,0) node[right]{time};
    \draw[-{Stealth[length=3pt]}] (0,-2.5) -- (0,2.7);
    \draw[dashed,intcol] (0,1.8) -- (8.5,1.8) node[right]{your limit};
    \draw[dashed,intcol] (0,-1.8) -- (8.5,-1.8) node[right]{my limit};
    \draw[thick,keycol] plot[smooth] coordinates {(0,0) (1,-0.8) (2,-1.5) (3,-0.6) (4,0) (5,1.2) (6,1.6) (7,0.5) (8,-0.4)};
    \node[left] at (0,0) {0};
    \node[font=\small\bfseries] at (4,3.0) {credit: net position};
  \end{scope}
\end{tikzpicture}
\caption{Payments with a fixed quantity of money (left: balances only move between pockets, and a trade is impossible when the buyer has
run out) and with credit (right: the net position starts at zero and expands and contracts within credit limits)
\ts{4}{4}{0:28}\ts{4}{4}{1:35}.}\label{fig:mv-de}
\end{nfigure}

This is where discipline and elasticity (\cref{sec:nh-discipline}) become concrete \ts{4}{4}{0:07}.
\begin{itemize}
  \item With a \textbf{fixed quantity of money}, my balance goes up when I sell to you and down when I buy, while the total does not
        change. If there is little money, people run out, and some mutually advantageous trades are impossible: that is
        \emph{discipline} \ts{4}{4}{0:28}\ts{4}{4}{0:49}\ts{4}{4}{1:13}.
  \item With \textbf{pure credit}, positions start at zero and expand and contract with the pattern of payments, within each party's credit
        limit; there is no money, only credit that comes and goes \ts{4}{4}{1:35}\ts{4}{4}{1:56}\ts{4}{4}{2:16}.
  \item With a \textbf{bank}, I can borrow up to my limit at the bank, and I am willing to hold large positive balances with the bank, much
        more than I would lend to you; the bank does this with many people at once: a multilateral credit system
        \ts{4}{4}{2:58}\ts{4}{4}{3:18}.
\end{itemize}

\begin{proposition}[The payment system is a credit system]\label{prop:mv-credit}
The payment system is a credit system, not a money system, and it has to be one in order to have enough elasticity: people who want to
make mutually advantageous trades must be able to make them even if none of them has money (gold, the thing fixed in quantity) at the
moment \ts{4}{4}{3:39}. When looking at the actual system, expect to find credit, even where it does not look like credit
\ts{4}{4}{4:00}.
\end{proposition}

% ------------------------------------------------------------------
\section{The survival constraint}\label{sec:mv-survival}

Mehrling reads from Minsky's \emph{Stabilizing an Unstable Economy} (1986), p.~198 \ts{4}{5}{0:07}: to analyse how financial commitments
affect the economy one must look at economic units in terms of their cash flows; the cash flow approach looks at all units, households,
corporations, state and municipal governments, even national governments, \emph{as if they were banks} \ts{4}{5}{0:28}. That is: every unit
is a money-flow operation, cash in and cash out, and the goods side is secondary \ts{4}{5}{0:49}. The macro side of this view comes from
Copeland (\cref{sec:mv-fof}); the micro side is the survival constraint \ts{4}{5}{1:09}\ts{4}{5}{1:30}.

\begin{definition}[New concepts: survival (reserve) constraint]\label{def:mv-survival}
The \emph{survival constraint} (Minsky), or \emph{reserve constraint}, of an agent is the requirement that in every period
\[
  \text{cash inflow} \;\ge\; \text{cash outflow (cash commitments)} ,
\]
where cash commitments include promises already made, such as debt repayments. If it fails, the agent defaults: ``it doesn't matter if
you have lots of wealth, if you have no cash'' \ts{4}{5}{2:10}\ts{4}{5}{2:31}\ts{4}{5}{2:51}.
\end{definition}

The binding constraint is not the budget constraint of economics, and certainly not an \emph{intertemporal} budget constraint: it is a
liquidity constraint \ts{4}{5}{1:50}\ts{4}{5}{2:10}. ``Liquidity kills you quick'': a bank can go on doing business for a long time while
insolvent, many do, but cannot survive one day illiquid. If everyone is a bank, this is what focuses the mind \ts{4}{5}{2:51}\ts{4}{5}{3:12}.

\begin{intuition}[Budget constraint versus survival constraint]
The intertemporal budget constraint, $\sum_t c_t/(1+r)^t \le W$ (\cref{eq:fp-budget}), restricts present values: any timing of payments
is allowed as long as the totals balance, because one can always borrow against future income at the rate $r$. The survival constraint is
a constraint at \emph{each} date: a pointwise condition rather than an integral one. Two agents with the same present values can
differ completely in whether they survive.
\end{intuition}

\paragraph{Three ways to meet a shortfall.} If cash inflow does not cover a commitment, there are three choices \ts{4}{9}{0:07}\ts{4}{9}{0:29}:
\begin{enumerate}
  \item \textbf{dishoard}: run down cash reserves;
  \item \textbf{borrow}: find someone to lend more cash;
  \item \textbf{sell assets}, at whatever price one can get: a fire sale.
\end{enumerate}
(One could also sell goods; the course focuses on financial assets \ts{4}{9}{0:49}.) Only the first is dependable in a crisis: the others
need someone to take the other side of the trade, at a good price, and in a crisis they may not want to \ts{4}{9}{1:09}\ts{4}{9}{1:31}. This is
why reserves matter: at the end of the day it is cash that meets commitments. But nobody wants to hold only money; credit lets us escape the
constraint by pushing it into the future: we live to fight another day, but that day the borrowing comes due and we face the constraint
again \ts{4}{9}{1:52}\ts{4}{9}{2:12}. \emph{Credit gives us elasticity today but discipline tomorrow} \ts{4}{9}{2:32}.

% ------------------------------------------------------------------
\section{Sources and uses of funds}\label{sec:mv-su}

The \emph{flow of funds} accounts are real accounts, collected by the Fed and watched closely by financial advisers; start from their
micro foundation \ts{4}{6}{0:06}\ts{4}{6}{0:27}. For one agent, distinguish \emph{uses} from \emph{sources} of funds, and the goods-and-services
account (``above the line'') from the financial account (``below the line''), which has three categories: financial assets, financial
liabilities and money, a hierarchy of credit and money in which money is the ultimate means of payment \ts{4}{6}{0:48}\ts{4}{6}{1:08}\ts{4}{6}{1:29}.

\begin{nfigure}
\renewcommand{\arraystretch}{1.35}
\begin{tabular}{l|c|c}
 & \textbf{Uses} (cash out) & \textbf{Sources} (cash in)\\\hline
Goods and services (incl.\ labour) & expenditure & receipts\\\hline
Financial assets & accumulation & decumulation\\
Financial liabilities & repayment & borrowing\\
Money & hoarding & dishoarding\\
\end{tabular}
\caption{The sources-and-uses framework \ts{4}{6}{1:53}\ts{4}{6}{4:01}. Every entry is a \emph{flow} over a period, not a stock: this is not a
T-account. Above the line, goods; below the line, the financial account.}\label{fig:mv-su}
\end{nfigure}

\begin{itemize}
  \item Selling anything, including labour, is a receipt, a source of funds; buying is an expenditure, a use. Unlike the national
        accounts, which count only final sales (value added), here every sale counts: it is a source of cash for the seller even if it
        adds no value for the economy \ts{4}{6}{2:14}\ts{4}{6}{2:38}.
  \item Selling a financial asset is a source (decumulation), buying one a use (accumulation) \ts{4}{6}{2:58}.
  \item Borrowing is a source, repaying a use \ts{4}{6}{3:19}.
  \item Taking cash out of one's pile is dishoarding (a source), adding to it hoarding (a use) \ts{4}{6}{3:41}\ts{4}{6}{4:01}.
\end{itemize}
The framework is complete: every transaction fits in one of the categories, which lets one sort out complicated transactions
\ts{4}{6}{4:22}.

\begin{proposition}[The two rules of the flow of funds]\label{prop:mv-rules}
\leavevmode
\begin{enumerate}[(i)]
  \item For each agent, every use has a corresponding source and vice versa: nothing can be bought, no asset accumulated, no debt
        repaid, no cash hoarded, without the money coming from somewhere \ts{4}{6}{4:42}\ts{4}{6}{5:03}\ts{4}{6}{5:24}.
  \item Every agent's use is some other agent's source and vice versa: my receipt is somebody's expenditure, my purchase of an asset is
        somebody's sale \ts{4}{6}{5:45}\ts{4}{6}{6:08}.
\end{enumerate}
Consequently each transaction requires at least four entries, two for each of two agents \ts{4}{7}{1:56}. Rule (ii) is what turns the
framework into macroeconomics: the economy is a web in which everything is connected \ts{4}{6}{6:28}.
\end{proposition}

\begin{intuition}[Double entry, twice]
Rule (i) is the conservation law of each agent's cash (inflow $=$ outflow, with hoarding as the ``storage'' term); rule (ii) is the
conservation law between agents (every outflow is someone's inflow). Together they make the matrix of all agents' flows have zero row sums
(each type of flow balances across agents) and equal column sums for each agent (uses $=$ sources), as Kirchhoff's laws do for currents at
nodes and around a circuit.
\end{intuition}

\subsection{Two examples}\label{sec:mv-examples}
\paragraph{A coffee paid in cash.} Before class Mehrling bought a \$2 coffee at Oren's Daily Roast, in cash \ts{4}{7}{0:08}\ts{4}{7}{0:28}:
\begin{center}
\begin{tabular}{l|cc|cc}
 & \multicolumn{2}{c|}{\textbf{Mehrling}} & \multicolumn{2}{c}{\textbf{Oren's}}\\
 & uses & sources & uses & sources\\\hline
goods & coffee 2 & & & coffee 2\\\hline
money & & dishoard 2 & hoard 2 & \\
\end{tabular}
\end{center}
Four entries, both rules satisfied \ts{4}{7}{0:48}\ts{4}{7}{1:09}\ts{4}{7}{1:35}. The point is to build intuition for the inside character of money,
for the interlocking of the monetary system, for its money flows; these intuitions are not natural for economists, trained to see money as a
veil \ts{4}{7}{1:56}\ts{4}{7}{2:17}.

\paragraph{A dinner paid by credit card.} After class, dinner at Virelli, \$20, charged to MasterCard \ts{4}{7}{2:38}\ts{4}{7}{3:00}. Three parties,
MasterCard in the middle as the bank of \cref{sec:mv-payments} \ts{4}{7}{3:25}. In stages \ts{4}{7}{3:46}--\ts{4}{7}{9:04}:
\begin{enumerate}
  \item \textbf{Dinner.} Mehrling's expenditure, Virelli's receipt. Signing the slip is not paying but \emph{borrowing}: an IOU to MasterCard
        (Mehrling borrows, MasterCard accumulates). Virelli lets him go because it gets MasterCard's IOU (MasterCard borrows, Virelli
        accumulates) \ts{4}{7}{4:33}\ts{4}{7}{4:53}\ts{4}{7}{5:14}\ts{4}{7}{5:36}. Two new credits, out of thin air.
  \item \textbf{Merchant settlement.} The next day or week Virelli presents its claim: Virelli decumulates and hoards cash, MasterCard dishoards
        and repays \ts{4}{7}{7:00}\ts{4}{7}{7:21}\ts{4}{7}{8:03}. For Virelli it is as if it had been paid in cash.
  \item \textbf{Cardholder settlement.} At the end of the month Mehrling repays (dishoarding cash, repaying debt), and MasterCard decumulates the
        IOU and hoards cash \ts{4}{7}{8:24}\ts{4}{7}{8:44}\ts{4}{7}{9:04}.
\end{enumerate}
\begin{center}
\begin{tabular}{l|cc|cc|cc}
 & \multicolumn{2}{c|}{\textbf{Mehrling}} & \multicolumn{2}{c|}{\textbf{MasterCard}} & \multicolumn{2}{c}{\textbf{Virelli}}\\
 & uses & sources & uses & sources & uses & sources\\\hline
goods & 1: dinner & & & & & 1: dinner\\\hline
assets & & & 1: IOU$_\text{M}$ & 3: IOU$_\text{M}$ & 1: IOU$_\text{MC}$ & 2: IOU$_\text{MC}$\\
liabilities & 3: repay & 1: IOU$_\text{M}$ & 2: repay & 1: IOU$_\text{MC}$ & & \\
money & & 3: dishoard & 3: hoard & 2: dishoard & 2: hoard & \\
\end{tabular}
\end{center}
Credit expanded, two credits at once, and then collapsed in two phases; at the end there is no credit left, and cash settled the debts: the
elasticity came from credit, the discipline from the debts having to be paid \ts{4}{7}{9:26}\ts{4}{7}{9:47}.

% ------------------------------------------------------------------
\section{The flow of funds accounts}\label{sec:mv-fof}

Morris Copeland, \emph{A Study of Moneyflows in the United States} (NBER, 1952), proposed these accounts as a framework for
macroeconomic analysis \ts{4}{8}{0:07}. They apply the sources-and-uses framework to whole sectors, with aggregation and netting, which is
what makes them hard to read for people who do not have the framework in mind \ts{4}{8}{0:27}.
\begin{itemize}
  \item \textbf{Columns}: sectors (households and nonprofit organizations, nonfinancial business, state and local governments, federal
        government, domestic financial sectors, rest of the world), each with a ``uses'' and a ``sources'' column \ts{4}{8}{0:48}\ts{4}{8}{2:31}.
  \item \textbf{Rows}: first the goods and services part (about 12 lines), then the financial instruments: various kinds of money first
        (official reserves, SDRs, Treasury currency, deposits; lines 14--22), then forms of credit (from line 23), in the reverse order
        of the hierarchy \ts{4}{8}{1:09}\ts{4}{8}{1:29}.
  \item \textbf{Checks}: by rule (i) each sector's uses and sources columns must have the same sum; by rule (ii) each row must sum to zero.
        They do not quite: the \emph{statistical discrepancies} show where the net has missed something \ts{4}{8}{1:50}\ts{4}{8}{2:10}.
  \item \textbf{Negative entries}: the published accounts record repayment as a negative source (a reduction of borrowing) and
        decumulation as a negative use; this moves entries to the other side but keeps the rules \ts{4}{8}{3:12}\ts{4}{8}{3:33}\ts{4}{8}{3:53}.
\end{itemize}

\paragraph{Copeland against Keynes and the monetarists} \ts{4}{8}{4:54}:
\begin{itemize}
  \item $C+I+G=Y$ ignores purchases and sales of financial assets and of existing goods (old houses, a very big market), and much of what
        happens in the monetary and financial system \ts{4}{8}{5:15}\ts{4}{8}{5:35};
  \item $MV=PT$ treats all transactions as made with money, ``as if the whole economy were Oren's Daily Roast'', while most transactions
        are credit transactions, absorbed into velocity \ts{4}{8}{5:56}.
\end{itemize}
Mehrling thinks Copeland was right, but he did not win: the accounts were compiled by a small division of the Board of Governors, but not
used for policy, because no macroeconomic theory was attached to them; Keynesians and monetarists each had theirs. ``So there's a Nobel
prize for you'' \ts{4}{8}{6:17}\ts{4}{8}{6:37}\ts{4}{8}{6:58}.

\paragraph{What the accounts miss.} In the table shown, the largest discrepancy in the last column is \$115 billion on line 22, ``Fed funds
and security RPs'': the wholesale money market escapes the statistical net \ts{4}{8}{7:18}\ts{4}{8}{7:38}. Sector discrepancies are large too
(\$561 billion for households and nonprofits, \$440 billion for domestic nonfinancial sectors), partly because ``nonfinancial'' firms are
also financial: GM makes cars, but it is also a kind of bank with international operations \ts{4}{8}{7:59}\ts{4}{8}{8:19}\ts{4}{8}{8:40}. And there
are no derivatives at all, no swaps, no options: they did not exist in 1952, and they are off balance sheet \ts{4}{8}{9:22}\ts{4}{8}{9:42}. After the
crisis the Office of Financial Research was created to build new accounting structures; Mehrling proposed to ``soup up'' the accounts to
include swaps and derivatives, as the course will do \ts{4}{8}{8:59}\ts{4}{8}{9:42}.\begin{coursenote}[Flow of funds today]
The flow of funds accounts have been published by the Fed since 1955, now as the \emph{Financial Accounts of the United States} (release
Z.1). The Office of Financial Research, created by the Dodd--Frank Act in 2010, is an office of the US Treasury Department, not of the Fed
(in class it is called ``a division of the Fed'' \ts{4}{8}{8:59}).
\end{coursenote}

% ------------------------------------------------------------------
\section{Liquidity, long and short}\label{sec:mv-longshort}

The survival constraint applies not only today but tomorrow and after: pushing a problem into the future lets you fight another day, but
you should push it where you expect cash inflows, not where you have other commitments \ts{4}{10}{0:28}\ts{4}{10}{0:48}. Example, over four
periods $t,\dots,t+3$, with cash inflow 10 every period (40 in all) and total commitments 20: no solvency problem, a ``profit'' of 20
\ts{4}{10}{1:31}\ts{4}{10}{1:53}\ts{4}{10}{2:14}.

\begin{nfigure}
\renewcommand{\arraystretch}{1.35}
\begin{tabular}{l|cccc|l}
 & $t$ & $t+1$ & $t+2$ & $t+3$ & \\\hline
cash inflow (all cases) & 10 & 10 & 10 & 10 & \\\hline
commitments, hedge & 5 & 5 & 5 & 5 & never binds\\
commitments, short-term & 20 & 0 & 0 & 0 & binds today\\
\end{tabular}
\caption{Same stocks, different flows \ts{4}{10}{2:35}\ts{4}{10}{3:48}: in both cases inflows total 40 and commitments 20 (ignoring
discounting), so the balance sheet is the same; only the time pattern differs.}\label{fig:mv-patterns}
\end{nfigure}

\begin{enumerate}
  \item \textbf{Hedge (liquid) financing}: commitments of 5 every period; the constraint never binds, the question ``how am I going to
        meet it?'' never arises \ts{4}{10}{2:35}\ts{4}{10}{2:56}.
  \item \textbf{Short-term financing}: the business is fine (10 every period) but financed with short-term debt, all due today: 20 owed,
        10 coming in. One of the three choices is forced now: use reserves, roll over the debt, or sell the business
        \ts{4}{10}{3:48}\ts{4}{10}{4:09}. Short-term borrowing is the sign of vulnerability, and real banks, with demand deposits and overnight
        repo funding, look like this all the time; hence the importance of the overnight market, enormous and invisible to households
        \ts{4}{10}{4:30}\ts{4}{10}{4:50}\ts{4}{10}{5:11}.
  \item \textbf{A problem coming later} (a third case, drawn on the board with ``a little leeway''): no problem today, but a future misalignment of flows and commitments that the agent, and its creditors,
        can see. While still liquid, the agent refinances, spreading the debt over time or prepaying part of it: \emph{anticipated} liquidity
        constraints have consequences for liquidity today, because tomorrow the markets may be frozen and a perfectly solvent agent could be
        ``dead'' \ts{4}{10}{6:00}\ts{4}{10}{6:21}\ts{4}{10}{6:41}\ts{4}{10}{7:02}.
\end{enumerate}

\paragraph{Price-insensitive borrowers.} The economy is a crowd of agents with different time patterns of flows and commitments, all trading
in the same market \ts{4}{10}{7:22}\ts{4}{10}{7:43}. An agent under stress must borrow now, whether the short rate is 10\% or 20\%: its demand
is price-inelastic, and the more such agents there are, the further rates can move, for reasons that have nothing to do with time preference
\ts{4}{10}{8:03}\ts{4}{10}{8:23}. Think of Bear Stearns or Lehman before their collapse: this is the stuff of financial crises, and all banks, with
short-term liabilities that can vanish tomorrow, look like this \ts{4}{10}{8:44}. Agents with surplus cash lend it into the market. Over the
business cycle the balance swings between easy and tight credit, and it shows up in the money market \ts{4}{10}{9:04}\ts{4}{10}{9:25}.

\begin{definition}[New concepts: hedge financing, rollover]
\emph{Hedge financing} (Minsky): commitments in every period are covered by expected cash inflows of the same period. To \emph{roll over}
a debt is to repay it with the proceeds of a new debt, pushing the commitment into the future. (Minsky also defines \emph{speculative}
units, whose near-term commitments exceed near-term inflows and must be rolled over, and \emph{Ponzi} units, which must borrow even to pay
interest.)
\end{definition}

% ------------------------------------------------------------------
\section{Financial fragility: flows and stocks}\label{sec:mv-fragility}

\begin{nfigure}
\begin{tikzpicture}[font=\small]
  \draw[thick] (0,0) -- (0,2.6);
  \fill (0,2.6) circle (2pt);
  \draw[thick] (-2.6,2.3) -- (2.6,2.9);
  \draw (-2.6,2.3) -- (-3.3,1.2) (-2.6,2.3) -- (-1.9,1.2);
  \draw[thick,intcol,fill=intcol!15] (-3.4,1.2) arc (180:360:0.75);
  \node[intcol,align=center,font=\scriptsize] at (-2.65,0.1) {cash\\flows};
  \draw (2.6,2.9) -- (1.9,1.8) (2.6,2.9) -- (3.3,1.8);
  \draw[thick,thmcol,fill=thmcol!15] (1.8,1.8) arc (180:360:0.75);
  \node[thmcol,align=center,font=\scriptsize] at (2.55,0.7) {cash\\commitments};
  \draw[-{Stealth[length=4pt]},very thick,keycol] (0,2.6) -- (0.9,3.6);
  \draw[keycol] (-1.2,3.5) arc (135:45:1.7);
  \node[keycol,font=\scriptsize,align=center] at (0,4.3) {money market rate of interest};
  \draw[thick] (-0.8,0) -- (0.8,0);
\end{tikzpicture}
\caption{The scales \ts{4}{11}{0:07}\ts{4}{11}{0:30}: the balance between the pattern of cash flows and the pattern of cash commitments, across the
whole economy, sets the money market interest rate.}\label{fig:mv-scales}
\end{nfigure}

The image, ``because we're in New York'', is the scales of Lady Liberty: one pan holds cash flows, the other cash commitments, and the
pointer is the money market interest rate \ts{4}{11}{0:07}\ts{4}{11}{0:30}\ts{4}{11}{0:50}. The Fed does something too, but behind the scenes there are
balance sheet facts: promises to pay at future dates that either are or are not lined up with the prospects of whoever made them; if they are
not, that person is forced to act. ``It's not about preference, it's about being forced to do something by the promises you've made in the
past'' \ts{4}{11}{1:10}\ts{4}{11}{1:31}. You buy a house, then lose your job: the cash flows go away, the commitments stay \ts{4}{11}{1:52}.

\begin{proposition}[Minsky's fragility, in the money view]\label{prop:mv-fragility}
Financial fragility is the balance of cash commitments (promises to pay) relative to the cash flows actually arising from activity
(wages, profits, sales). When most agents' cash flows cover their commitments there is no crisis; when many agents' do not, the system
becomes fragile, and the crisis shows up in the money market rate of interest \ts{4}{11}{2:12}\ts{4}{11}{2:33}\ts{4}{11}{2:53}.
\end{proposition}

The money market is thus the ``crisis detector'', the ``stress detector'': the bouncing of its rate reflects the shifting balance between cash
flows and commitments, a ``sufficient statistic'' for the state of the macroeconomic balance \ts{4}{11}{3:15}\ts{4}{11}{3:36}.

\paragraph{Stocks are residues of flows.} In a T-account, financial assets and money on the left, financial liabilities and net worth on the
right; net worth is solvency \ts{4}{11}{3:57}\ts{4}{11}{4:17}. Each stock is the accumulation of past flows: assets bought, cash hoarded, debts
incurred \ts{4}{11}{4:38}\ts{4}{11}{4:58}. And each looks forward too: liabilities and assets have dates, the dates of future payments; the stock
value adds them up, as a discounted present value (40 and 20 in the examples) \ts{4}{11}{5:18}. The examples of \cref{fig:mv-patterns} have
identical stocks and very different prospects: the balance sheet shows solvency; to see fragility one must look behind it, at the time
patterns \ts{4}{11}{5:38}\ts{4}{11}{5:59}.

% ------------------------------------------------------------------
\flushside
\section*{Key formulas and ideas}
\begin{keyformulas}
\begin{tabularx}{\linewidth}{@{}l L@{}}
Money vs credit payments & money: existing balances move, total fixed; credit: outstanding IOUs increase with the transaction\\
Payment system & a credit system (elasticity), settled in money (discipline)\\
Survival constraint & every period: cash inflow $\ge$ cash commitments (not an intertemporal budget constraint)\\
Shortfall & dishoard, borrow, or sell assets; only reserves are dependable in a crisis\\
Sources and uses & receipts/expenditure; decumulation/accumulation; borrowing/repayment; dishoarding/hoarding\\
Two rules & (i) each agent: uses $=$ sources; (ii) each use is another agent's source $\Rightarrow$ $\ge4$ entries per transaction\\
Fragility & commitments vs cash flows in time; shows up in the money market rate\\
Stocks & accumulated past flows and discounted future flows; same stocks can hide different time patterns\\
\end{tabularx}
\end{keyformulas}

\section*{Exercises}

\begin{exercise}[Sources and uses]
Record in the sources-and-uses framework, for all agents involved:
\begin{enumerate}[(a)]
  \item you receive a \$3000 salary paid into your bank account;
  \item you buy \$1000 of shares from another investor, paying from your account;
  \item you repay \$500 of a bank loan;
  \item a firm sells a used machine to another firm on 30-day trade credit.
\end{enumerate}
Check both rules in each case.
\end{exercise}

\begin{exercise}[The credit card, completed]
Write the three stages of the dinner example as changes in T-accounts (stocks) for Mehrling, MasterCard and Virelli, assuming each pays
through a bank. What is the total quantity of credit outstanding after each stage?
\end{exercise}

\begin{exercise}[Time patterns]
A firm expects cash inflows of 10, 10, 10, 10 and has debt of 24.
\begin{enumerate}[(a)]
  \item Give a repayment schedule that never violates the survival constraint and one that violates it at $t$.
  \item In the second case, how much must it roll over at $t$? At what rate does it become insolvent if it must pay interest $i$ on the
        rolled-over amount for three periods?
  \item Explain why its demand for funds at $t$ is price-inelastic.
\end{enumerate}
\end{exercise}

\begin{exercise}[Statistical discrepancy]
In a two-sector economy households lend 100 to firms through repo, and firms report borrowing 85. Explain how the discrepancy shows up in the
row and column checks of the flow of funds matrix, and where you would look for the missing flow.
\end{exercise}
